\documentclass[10pt,a4paper]{amsart}
\usepackage[margin=19mm]{geometry}
\usepackage{amsmath,amssymb,mathtools}
\usepackage[scr=rsfs,cal=euler]{mathalfa}
\usepackage{xfrac}
\usepackage[unicode,hidelinks,bookmarks=false]{hyperref}
\newcommand{\ie}{i.e.\ }
\newcommand{\cf}{cf.\ }
\newcommand{\resp}{respectively\ }
\newcommand{\Subsection}{\subsection}
\providecommand{\nobreakdash}{\nobreak\hbox{-}}
\setlength{\emergencystretch}{2em}
\multlinegap0pt
\usepackage{bm}
\usepackage{bbm}
\usepackage{dsfont}
\usepackage{xspace}
\usepackage{cancel}
\usepackage{mathtools}
\usepackage{amsthm}
\makeatletter
\@ifundefined{lemma}{\newtheorem{lemma}{Lemma}}{}
\@ifundefined{theorem}{\newtheorem{theorem}{Theorem}}{}
\makeatother
\title[A note on solitary numbers]{A note on solitary numbers}
\author{Sagar Mandal}
\date{Source: arXiv:2503.11694v2 (CC BY 4.0)}
\begin{document}
\maketitle

\noindent\emph{Source and licence.} A note on solitary numbers. Version arXiv:2503.11694v2 (CC BY 4.0), distributed under the Creative Commons Attribution 4.0 International licence, as recorded in the arXiv licence metadata for this exact version and shown on the abstract page https://arxiv.org/abs/2503.11694v2. Full rights notice: licenses/arxiv-2503.11694-CC-BY-4.0.txt.\par\smallskip

\noindent\textit{Editorial scope.} Counted unit: the Theorem whose source label is \texttt{thm 3.1} in the article (source0.tex lines 125--127, source label \texttt{thm 3.1}), delivered together with the article's own complete original proof of it (source0.tex lines 128--150, one \texttt{proof} environment of 23 lines). No mathematics is written, repaired or re-derived: statement and proof are the article's own text line for line. Required context retained uncounted: Lemma 2.1 (lines 83--85); Lemma 2.2 (lines 88--90); Lemma 2.5 (lines 113--115). Every definition, display and label of those lines is retained unchanged. Omitted from this package and not redistributed: 1--82, 86--87, 91--112, 116--124, 151--342. Nothing inside a retained excerpt is omitted or altered: every retained body is delivered exactly as the article's lines are written, so no omission receipt of any kind accompanies this package. Citations: the retained excerpts carry no citation command at all, so no bibliography entry of the article is transcribed and no reference is redistributed. Reference adaptation: none was needed and none was made; every reference inside the delivered unit resolves inside it, so no unresolved reference or citation is produced. Printed slips kept literal: no printed word, symbol, hypothesis or number of the counted statement or of its proof was altered. Layout and class: the article's own class is \texttt{amsart} with packages including tabularx, geometry, amssymb, amsmath, amsthm, latexsym, mathrsfs, url, amsfonts, graphicx, hyperref, mathtools, enumerate, xcolor; the wrapper is the lane's pinned \texttt{amsart} editorial wrapper at 10pt on A4, which restates the theorem-like environments the retained text needs and adds the article's own macro definitions that the retained text needs (none), with extra wrapper packages (bm, bbm, dsfont, xspace, cancel, mathtools). The delivered numbering is the unit's own. Assets: no retained span contains a figure environment, a graphics-inclusion command, a PDF-page inclusion, a table or a TikZ picture; no image file is copied into this package, the package declares no asset dependency and no image file is redistributed with it. Licence: version arXiv:2503.11694v2 of this article is distributed under the Creative Commons Attribution 4.0 International licence, as recorded in the arXiv licence metadata for this exact version and shown on the abstract page https://arxiv.org/abs/2503.11694v2. A full rights notice is delivered at licenses/arxiv-2503.11694-CC-BY-4.0.txt. Characters: every retained excerpt is pure ASCII, so the delivered engine, XeLaTeX with its default Latin Modern text font, reports no missing character.
\begin{lemma}\label{Lemma 2.1}
$I(n)$ is weakly multiplicative, that is, for any two co-prime positive integers $n$ and $m$ we have $I(nm)=I(n)I(m)$.
  \end{lemma}

\begin{lemma}\label{Lemma 2.2}
If $\gamma,n$ are two positive integers and $\gamma>1$. Then $I(\gamma n)>I(n)$.
\end{lemma}

\begin{lemma}\label{Lemma 2.5}
If $n=\displaystyle{\prod_{i=1}^{k}p_i^{\gamma_i}}$, then $\displaystyle{I(n)<\prod_{i=1}^{k}\frac{p_i}{p_i-1}}$.
\end{lemma}

\begin{theorem}\label{thm 3.1}
Let $F$ be a friend of $14$, then $F$ is an odd positive non-square integer.   
\end{theorem}

\begin{proof}
    Let a positive integer $F$ be a friend of $14$. Then $F$ must be greater than $14$ as for any positive integer less than $14$ we have $I(F)\neq \frac{12}{7}$. Since $\frac{\sigma(F)}{F}=I(F)=\frac{12}{7}$ we have $7\sigma(F)=12F$, from which it follows that $7\mid F$ and $12\mid \sigma(F)$ as $(7,12)=1$. Therefore, we can write $F=7F'$ for some positive integer $F'>2$. 

Let us  assume that $F'$ is an even positive integer. Then we can rewrite $F$ as $F=14F''$ for some positive integer $F''>1$, but then $I(F)>I(14)$ by Lemma \ref{Lemma 2.2}. Therefore, $F'$ is an odd positive integer and thus $F$ is an odd positive integer.

    Now if $F=7^a$ for some positive integer $a>2$, then it cannot be a friend of $14$ as by Lemma~\ref{Lemma 2.5} we have $I(7^{a})<\frac{7}{6}<\frac{12}{7}$. Therefore, $F$ must be written as $\displaystyle{F=7^{a_{1}}\cdot \prod_{i=2}^{\omega(F)}p_{i}^{a_{i}}}$($p_1=7$). 

Let us suppose that all $a_i$ are even. Then
$$I(F)=\frac{\sigma(F)}{F}=\frac{12}{7},$$
using Lemma \ref{Lemma 2.1} we get
\begin{align*}
   I(7^{a_1})\cdot \prod_{i=2}^{\omega(F)}I(p_{i}^{a_i})=\frac{12}{7}.
   \end{align*}
   This implies
   \begin{align*}
 \sigma(7^{a_1})\cdot \prod_{i=2}^{\omega(F)}\sigma(p_i^{a_i})=12\cdot7^{{a_1}-1}\cdot\prod_{i=2}^{\omega(F)}p_{i}^{a_i},
\end{align*}
that is
\begin{align*}
 (1+7+ \cdots +7^{a_1})\cdot\prod_{i=2}^{\omega(F)}(1+p_i+\dots+p_i^{a_i})=12\cdot7^{{a_1}-1}\cdot\prod_{i=2}^{\omega(F)}p_{i}^{a_i}.
\end{align*}
Since $p_i>2$, for all $1\leq i\leq \omega(F)$ the right-hand side of the above expression is an even integer but the left-hand side is odd since $a_i$ are even, which immediately implies that $(1+p_i+\dots+p_i^{a_i})$ is odd, which is absurd. Hence, all $a_i$ cannot be even integers. Therefore, $F$ is a non-square positive integer. This proves that $F$ is an odd positive non-square integer.
\end{proof}

\end{document}
